\abstract{РЕФЕРАТ}

Объем работы равен \formbytotal{lastpage}{страниц}{е}{ам}{ам}. Работа содержит \formbytotal{figurecnt}{иллюстраци}{ю}{и}{й}, \formbytotal{tablecnt}{таблиц}{у}{ы}{}, \arabic{bibcount} библиографических источников и \formbytotal{числоПлакатов}{лист}{}{а}{ов} графического материала. Количество приложений – 2. Фрагменты исходного кода представлены в приложении А. Графический материал представлен в приложении Б.

Перечень ключевых слов: оперативное перепланировение, производственное планирование, диспетчеризация, гибкая производственная система, нейросетевые модели, обучение с подкреплением, поведенческое клонирование, самоконтролируемое обучение, эвристические алгоритмы, оптимизация расписаний, цифровой двойник производства.

Объектом исследования являются процессы оперативного производственного планирования на промышленных предприятиях с дискретным характером производства.

Целью выпускной квалификационной работы является разработка интеллектуальной системы, обеспечивающей повышение эффективности производственного планирования посредством адаптивной диспетчеризации операций с использованием нейро-сетевых методов.

При проведении исследований использовались методы системного анализа, математического моделирования, эвристической оптимизации, глубокого обучения и объектно-ориентированного программирования.

Разработана интеллектуальная система производственного планирования с графическим интерфейсом, включающая модули управления заказами и ресурсами, построения расписаний и обучения нейросетевых диспетчеров.

При создании программного продукта использовалась двухзвенная архитектура, технология объектно-ориентированного программирования и библиотеки глубокого обучения PyTorch. Для визуализации результатов применялись диаграмма Ганта и средства аналитической графики Matplotlib.

Программный продукт может быть интегрирован в корпоративную информационную среду предприятия и использоваться в качестве модуля оперативного планирования и диспетчеризации.

\selectlanguage{english}
\abstract{ABSTRACT}
  
The volume of work is \formbytotal{lastpage}{page}{}{s}{s}. The work contains \formbytotal{figurecnt}{illustration}{}{s}{s}, \formbytotal{tablecnt}{table}{}{s}{s}, \arabic{bibcount} bibliographic sources and \formbytotal{числоПлакатов}{sheet}{}{s}{s} of graphic material. The number of applications is 2. The graphic material is presented in annex A. The layout of the site, including the connection of components, is presented in annex B.

List of keywords: operational rescheduling, production planning, dispatch, flexible production system, neural network models, reinforcement learning, behavioral cloning, self-supervised learning, heuristic algorithms, schedule optimization, digital production twin.
The object of the study is the processes of operational production planning at industrial enterprises with a discrete nature of production.

The purpose of the final qualifying work is to develop an intelligent production planning system, ensuring the construction of executable schedules and adaptive dispatching of operations based on neural network methods.

When conducting research, methods of system analysis, mathematical modeling, heuristic optimization, deep learning  and object-based learning were used.
oriented programming.

An intelligent production planning system with a graphical interface has been developed, including modules for managing orders and resources, creating schedules and training neural network dispatchers.

When creating the software product, a two-tier architecture was used, object-oriented programming technology and deep learning libraries PyTorch. To visualize the results, a Gantt chart and Matplotlib analytical graphics tools were used.

The software product can be integrated into the corporate information environment of an enterprise and used as a module for operational planning and dispatch.
\selectlanguage{russian}
